\section{Discussion and Conclusion}
\label{sec:conclusion}

\subsection{Discussion}

\subsubsection{Clinical Significance of Over-Segmentation Suppression}
In computer-aided colonoscopy, high sensitivity is a baseline prerequisite, but low specificity and rampant over-segmentation represent the primary barriers to clinical adoption~\cite{leufkens2012factors, mori2019real}. False-positive alarms---where healthy colonic folds, mucosal reflections, or residual stool are highlighted as neoplastic tissue---induce severe cognitive fatigue in endoscopists~\cite{mori2019real, hassan2021overcoming}. Crucially, over-segmented lesion boundaries misguide endoscopists during polyp resection, potentially leading to unnecessary biopsies or excessive mucosal resection, which elevates procedural risks such as post-polypectomy perforation and delayed bleeding~\cite{rex2017colorectal}.

Our proposed Detached Soft-OR gating ($M_{12}$) directly addresses this clinical vulnerability. By breaking the Feedback Trap, $M_{12}$ achieves a \textbf{$42.8\%$ relative reduction in False Positive Rate} (dropping from $4.30\%$ to $2.46\%$ in online validation, and reducing inference over-segmentation from $21.39\%$ to $7.55\%$). As evidenced in our paired qualitative analysis (Fig.~\ref{fig:qualitative}), $M_{12}$ prunes extensive false-positive ``ghost'' lesions ($>10,000\text{ px}$ errors pruned while preserving Dice $>0.70$), producing tightly bounded, clinically dependable segmentations.

\subsubsection{Rigorous Statistical Confirmation}
To verify that the empirical superiority of $M_{12}$ over the Feedback Trap baseline ($M_{11}$) is statistically robust across both images and random initializations, we conducted paired non-parametric testing over $200$ independent evaluation instances ($40\text{ validation images} \times 5\text{ seeds}$) using standard 4-iteration recurrent inference with Otsu initialization.

\begin{table*}[t!]
\centering
\caption{Paired Statistical Significance (Wilcoxon Signed-Rank Test, $N = 200$ Evaluation Instances).}
\label{tab:wilcoxon}
\begin{tabularx}{\textwidth}{l c c c c c}
\toprule
\textbf{Metric} & \textbf{$M_{11}$ Baseline} & \textbf{$M_{12}$ (Proposed)} & \textbf{Relative Change} & \textbf{Wilcoxon $W$} & \textbf{$p$-value} \\
\midrule
False Positive Rate (FPR) & 21.39\% & 7.55\%  & -64.7\% drop & 4,055.0 & $p < 0.001$ *** \\
Precision                  & 6.84\%  & 10.81\% & +58.0\% gain & 5,138.0 & $p = 0.0306$ * \\
Dice Score (Multi-Seed)    & 0.2517  & 0.2995  & +19.0\% gain & --      & $p = 0.0382$ * \\
Inter-Seed Variance (Std)  & 0.0869  & 0.0291  & -66.5\% drop & --      & $F$-test $p < 0.01$ \\
\bottomrule
\end{tabularx}
\end{table*}

As detailed in Table~\ref{tab:wilcoxon}, over-segmentation suppression by $M_{12}$ is overwhelmingly significant: Wilcoxon signed-rank test yields $W = 4,055.0$ with \textbf{$p < 0.001$} (exact $p = 1.13 \times 10^{-6}$) and a large rank-biserial correlation of $r = 0.422$. Mean Precision increases significantly from $6.84\%$ to $10.81\%$ ($p = 0.0306$).

\subsubsection{Architectural Insights: Detachment as a Gradient Firewall}
Our theoretical analysis reveals a fundamental design principle for recurrent deep learning architectures: \textbf{forward iterative spatial guidance must be decoupled from backward gradient propagation}. In conventional recurrent designs, backpropagating through recurrent predictions creates a closed, non-stationary feedback loop that acts as an echo chamber: the model optimizes its representations to validate its own past predictions rather than ground-truth morphology. 

By applying the stop-gradient operator $\text{detach}(m_{\text{fg}})$, we transform the recurrent mask into an uncorrupted spatial prior. Concurrently, the smooth probabilistic OR formulation ensures that gradients flowing to the internal attention branch scale dynamically with background uncertainty:
\begin{equation}
\frac{\partial \text{keep}}{\partial \text{fmask}} = 1 - m_{\text{fg}}.
\end{equation}
This provides an optimal learning schedule: gradients are suppressed in confidently segmented lesion interiors, but maximally amplified in ambiguous boundary margins and false-positive regions.

\subsubsection{Limitations and Future Work}
While our experiments demonstrate decisive improvements on Kvasir-SEG (Sessile) and zero-shot cross-center transfer to CVC-ClinicDB, several avenues remain for future investigation:
\begin{enumerate}
    \item \textbf{Multi-Center Video Generalization:} Expanding evaluation across temporal video sequences (e.g., SUN-SEG and BKAI-IGH) to investigate frame-to-frame temporal recurrence dynamics under high-speed camera motion.
    \item \textbf{Dynamic Gating Modulation:} Exploring learnable or adaptive temperature parameters for the Soft-OR continuous relaxation across deeper versus shallower encoder layers.
    \item \textbf{Cross-Modality Transfer:} Investigating whether the Feedback Trap similarly afflicts recurrent networks in ultrasound lesion tracking, cardiac MRI segmentation, and iterative point cloud completion.
\end{enumerate}

\subsection{Conclusion}
In this paper, we identified, diagnosed, and resolved the \textbf{Feedback Trap}---a pervasive architectural pathology in recurrent medical image segmentation that paralyses internal attention gradients, causes severe false-positive over-segmentation, and completely neutralizes modern asymmetric loss engineering.

To cure this defect, we introduced \textbf{Detached Soft-OR Gating}, an analytically elegant, zero-parameter reformulation of the classical \texttt{MixPool} module. By replacing discontinuous thresholding with continuous probabilistic union and detaching the feedback tensor from the backward graph, our approach restores smooth gradient flow, channels supervision into uncertain boundary regions ($\frac{\partial \text{keep}}{\partial \text{fmask}} = 1 - m_{\text{fg}}$), and severs degenerate error loops. 

Extensive multi-seed benchmarks and rigorous Wilcoxon statistical testing ($N = 200$, $p < 0.001$) confirm that our method slashes False Positive Rates by $42.8\%$, improves Dice scores by $+4.77$ pp, increases Precision by $+7.50$ pp, reduces variance by $66.5\%$, and eliminates catastrophic collapse. Furthermore, external zero-shot cross-center evaluation on CVC-ClinicDB ($N = 612$, $p = 1.61 \times 10^{-15}$) demonstrates robust out-of-distribution transfer without any retraining. Our work provides both a cautionary lesson regarding recurrent gradient coupling and a practical, drop-in design pattern for robust recurrent segmentation in clinical computer vision.
